% Source: https://artofproblemsolving.com/community/c3h1293339_candle_problem
% Stable identity: aops-topic-1293339-q01
After $t$ hours the heights (as fractions of the original) are $1-\dfrac {t}{4}$ and $1-\dfrac {t}{5}$; the second is taller. We need $1-\dfrac {t}{5}=2\left(1-\dfrac {t}{4}\right)$, i.e. $\dfrac {t}{2}-\dfrac {t}{5}=1$, so $t=\dfrac{10}{3}$. After $\boxed{3\text{ hours }20\text{ minutes}}$ the slower candle is twice as tall as the faster one ($\dfrac{1}{3}$ versus $\dfrac{1}{6}$ of the original height).
